\section{Related Work and Expected Contribution}
Recent research in single-cell transcriptomics has focused on both computational clustering and automated biological interpretation. Pan et al. introduced scMINER in 2025 as a mutual-information-based framework that combines cell clustering with transcription-factor and signaling-network inference \cite{pan2025scminer}. Its clustering component, MICA, uses mutual information to capture nonlinear relationships between cells and was evaluated against several commonly used single-cell clustering methods. In addition to identifying cell populations, scMINER provides marker genes and regulatory-driver information that can support downstream biological analysis.

Large language models have also been increasingly studied for the interpretation of single-cell data. Hou and Ji evaluated GPT-4 for cell-type annotation using marker-gene lists and showed that language models can assist in assigning biological identities to cell clusters \cite{hou2024gpt4}. More recently, Crowley et al. introduced AnnDictionary and benchmarked several large language models for de novo cell-type annotation using marker genes derived from unsupervised clustering \cite{crowley2025anndictionary}. Their results showed that LLM-based annotation can be useful, although performance varies across models, datasets, and cell populations.

More recent studies have incorporated retrieval-based methods to provide external biological evidence during cell-type annotation. Kazmi et al. explored a retrieval-augmented generation approach in which marker-gene information was supplemented with relevant biomedical literature before generating cell-type annotations \cite{kazmi2025beyond}. Yang et al. similarly developed a RAG-based framework that combined lexical and semantic retrieval, query expansion, reranking, and ontology-based evaluation for single-cell annotation \cite{yang2026rag}. These studies indicate that incorporating retrieved biomedical evidence can improve the grounding and reliability of language-model-generated annotations.

\subsection{Closest work and how this project differs}
Three groups of papers sit closest to what we propose. On scale, minicore clusters several million cells with information-theoretic divergences \cite{baker2021minicore}, and GPU pipelines now run standard single-cell workflows on atlas-sized data \cite{dicks2026rapids}. Neither computes mutual information between cells, which is the quantity MICA relies on and the one that has never been scaled. On choosing partitions without labels, sc-SHC tests each split for significance \cite{grabski2023significance} and CHOIR merges clusters that a classifier cannot separate \cite{sant2025choir}. Both assume a standard distance pipeline and neither has been applied to mutual information clustering. On interpretation, agent-style annotators such as CASSIA \cite{xie2025cassia} and methods that let a language model pick the clustering resolution \cite{song2025gptanno} still take marker genes as their input, while regulator activity is computed by tools like SCENIC \cite{aibar2017scenic} and scMINER but never handed to a retriever. Our project sits in the space these three groups leave open.

\subsection{Expected contribution}
If the project succeeds, we expect four contributions.

\begin{itemize}
\item An open Python implementation that makes exact mutual information neighbor graphs practical at a few million cells on one workstation, with measured neighbor recall and a reported fraction of neighbors certified exact.
\item The first accuracy evaluation of mutual information clustering at atlas scale and without using labels to pick the number of clusters, including how well small populations survive.
\item A measurement of what the cell-to-cell mutual information actually captures, separating co-detection from expression level and testing sensitivity to sequencing depth, together with simple estimator fixes if the test shows a problem.
\item The first test of whether regulatory driver activity works as retrieval evidence for cluster interpretation, measured against marker genes alone.
\end{itemize}

We want to be honest about one outcome. It is possible that mutual information clustering does not beat standard pipelines once labels are removed or once the data reaches atlas scale. That result would still be worth reporting, because nobody has measured it, and the tooling and evaluation protocol would stand on their own.
